\documentclass[11pt]{article}

\usepackage[T1]{fontenc}
\usepackage[utf8]{inputenc}
\usepackage{lmodern}
\usepackage[margin=1in]{geometry}
\usepackage{amsmath,amssymb,amsthm,mathtools}
\usepackage{microtype}
\usepackage{xcolor}
\usepackage{booktabs,tabularx,array}
\usepackage{enumitem}
\usepackage[hidelinks]{hyperref}
\usepackage[nameinlink,noabbrev]{cleveref}

\definecolor{deepblue}{RGB}{24,70,120}
\definecolor{softblue}{RGB}{238,245,252}
\definecolor{softred}{RGB}{253,242,241}
\definecolor{softgreen}{RGB}{239,248,241}
\definecolor{darkred}{RGB}{145,36,36}
\definecolor{darkgreen}{RGB}{33,105,56}

\newtheorem{theorem}{Theorem}[section]
\newtheorem{proposition}[theorem]{Proposition}
\newtheorem{lemma}[theorem]{Lemma}
\newtheorem{corollary}[theorem]{Corollary}
\theoremstyle{definition}
\newtheorem{definition}[theorem]{Definition}
\newtheorem{question}[theorem]{Question}
\theoremstyle{remark}
\newtheorem{remark}[theorem]{Remark}

\newcommand{\ellTwo}{\ell_2}
\newcommand{\Sub}{\operatorname{Sub}}
\newcommand{\Spec}{\operatorname{Spec}}
\newcommand{\iso}{\cong}
\newcommand{\embeds}{\hookrightarrow}
\newcommand{\N}{\mathbb N}
\newcommand{\Ezero}{E_0}

\setlist[itemize]{leftmargin=1.5em,itemsep=0.25em,topsep=0.35em}
\setlist[enumerate]{leftmargin=1.7em,itemsep=0.3em,topsep=0.35em}

\newcommand{\statusbox}[2]{%
  \begin{center}
  \fcolorbox{#1}{#2}{\begin{minipage}{0.91\textwidth}\vspace{0.5em}
  \textbf{Status.} The submitted argument does not prove its theorem. The theorem itself is correct, by a known result of Anisca in the now-standard generality recorded in \cite[Theorem 6.6]{CCRF2023}. The full Godefroy problem remains open in the literature located through 20 June 2026.
  \vspace{0.5em}\end{minipage}}
  \end{center}}

\title{\textbf{Godefroy's Problem 7}\\[0.3em]
\large Audit of a partial argument, a corrected theorem, and the current obstruction}
\author{}
\date{Literature checked through 20 June 2026}

\begin{document}
\maketitle
\vspace{-2.0em}
\statusbox{deepblue}{softblue}

\begin{abstract}
Godefroy's question asks whether every separable non-Hilbertian Banach space contains a countably infinite family of pairwise non-isomorphic closed infinite-dimensional subspaces.  The submitted note attempts to prove the assertion when the ambient space is not \(\ellTwo\)-saturated.  Its central minimal-class step is invalid: minimality for isomorphic embeddability gives mutual embeddability, not isomorphism.  This distinction is essential because the Banach-space Schroeder--Bernstein principle fails, even for complemented embeddings.

The claimed partial theorem is nevertheless true.  A short correct proof follows from the theorem that a separable Banach space with only finitely many subspace isomorphism types contains \(\ellTwo\).  Consequently, every separable Banach space that is not \(\ellTwo\)-saturated has infinitely many pairwise non-isomorphic subspaces.  This result is known, not a new solution.  The unresolved case is precisely the non-Hilbertian \(\ellTwo\)-saturated regime.  We record the strongest structural restrictions presently available and explain why several natural attempts stop at open problems concerning \(d_2\)-minimal, \(d_2\)-hereditarily indecomposable, unconditional, and twisted Hilbert spaces.
\end{abstract}

\section{The problem and the verdict}

The source survey states the following problem \cite[Problem 7]{DLT2010}.

\begin{question}[Godefroy]
Let \(X\) be a separable non-Hilbertian Banach space.  Must there exist a sequence \((X_n)_{n\in\N}\) of closed infinite-dimensional subspaces of \(X\) that are pairwise non-isomorphic?
\end{question}

For a separable Banach space \(X\), write
\[
  \Spec(X)=\{[Y]_{\iso}:Y\in\Sub(X)\},
\]
where \(\Sub(X)\) denotes the closed infinite-dimensional subspaces of \(X\).  In the usual ZFC setting, failure of Godefroy's conclusion is equivalent to \(\Spec(X)\) being finite: every infinite set contains a countably infinite subset.

\begin{center}
\begin{tabularx}{0.95\textwidth}{>{\bfseries}p{0.24\textwidth}X}
\toprule
Item & Verdict \\
\midrule
The partial theorem & Correct, but already known. \\
The submitted proof & Incorrect at the minimal-preorder step. \\
A valid replacement & Apply the finite-spectrum theorem of Anisca, as recorded in \cite[Theorem 6.6]{CCRF2023}. \\
The full question & Still open in the sources located through 20 June 2026. \\
Remaining regime & Separable, non-Hilbertian, \(\ellTwo\)-saturated, near-Hilbert spaces with very strong minimality properties. \\
\bottomrule
\end{tabularx}
\end{center}

\section{The exact gap in the submitted argument}

The note orders subspace isomorphism classes by
\[
  A\preceq B
  \quad\Longleftrightarrow\quad
  \text{a representative of }A\text{ embeds isomorphically into a representative of }B.
\]
This is a preorder, not a partial order.  Antisymmetry fails in Banach-space theory.

\begin{proposition}[What minimality actually gives]
Let \(A_0\) be minimal in a finite preorder and suppose \(B\preceq A_0\).  Then minimality yields \(A_0\preceq B\), but it need not yield \(B=A_0\).
\end{proposition}

\begin{proof}
If \(A_0\not\preceq B\), then \(B\preceq A_0\) is a strict descent, contradicting minimality.  Hence \(A_0\preceq B\).  Equality is an additional antisymmetry conclusion and does not follow in a preorder.
\end{proof}

Applied to a subspace \(Z\subseteq Y\), the argument therefore proves only
\[
  Z\embeds Y
  \quad\text{and}\quad
  Y\embeds Z,
\]
where the first embedding is inclusion and the second comes from minimality.  It does \emph{not} prove \(Y\iso Z\).  Thus it produces a minimal space (one that embeds into each of its subspaces), not a homogeneous space (one that is isomorphic to each of its subspaces).

This is not a removable technicality.  Gowers constructed non-isomorphic Banach spaces each of which is isomorphic to a complemented subspace of the other \cite{GowersSB}.  Hence even mutual \emph{complemented} embeddability does not imply isomorphism.  The homogeneous-space theorem cannot be invoked at this point.

\begin{center}
\fcolorbox{darkred}{softred}{\begin{minipage}{0.90\textwidth}
\textbf{Invalid inference in the submitted proof:}
\[
  [Z]\preceq [Y]\text{ and }[Y]\text{ minimal}
  \quad\not\Longrightarrow\quad [Z]=[Y].
\]
The valid conclusion is only that \([Z]\) and \([Y]\) lie in the same mutual-embeddability component.
\end{minipage}}
\end{center}

\section{A correct proof of the partial theorem}

A Banach space \(X\) is \(\ellTwo\)-saturated if every member of \(\Sub(X)\) contains a further subspace isomorphic to \(\ellTwo\).

The key known input is the following theorem, attributed to Anisca and stated in this generality in \cite[Theorem 6.6]{CCRF2023}.  Anisca's original construction is in \cite{Anisca2007}; the later near-Hilbert theorem removes the finite-cotype restriction in the form needed here \cite{Cuellar2018}.

\begin{theorem}[Finite subspace spectrum forces \(\ellTwo\)]\label{thm:anisca}
If a separable Banach space has only finitely many isomorphism classes of closed infinite-dimensional subspaces, then it contains an isomorphic copy of \(\ellTwo\).
\end{theorem}

\begin{theorem}[Corrected non-\(\ellTwo\)-saturated result]\label{thm:corrected}
Let \(X\) be a separable infinite-dimensional Banach space.  If \(X\) is not \(\ellTwo\)-saturated, then \(X\) contains a countably infinite family of pairwise non-isomorphic closed infinite-dimensional subspaces.
\end{theorem}

\begin{proof}
Because \(X\) is not \(\ellTwo\)-saturated, there exists \(W\in\Sub(X)\) containing no copy of \(\ellTwo\).  If \(\Spec(W)\) were finite, \cref{thm:anisca} would imply that \(W\) contains \(\ellTwo\), a contradiction.  Thus \(\Spec(W)\) is infinite.  Choose one representative from each of countably many distinct classes in \(\Spec(W)\).  These representatives are subspaces of \(X\) and are pairwise non-isomorphic.
\end{proof}

\begin{corollary}\label{cor:saturated}
If a separable Banach space has no countably infinite pairwise non-isomorphic family of subspaces, then it is \(\ellTwo\)-saturated.  Therefore every counterexample to Godefroy's problem must be separable, non-Hilbertian, and \(\ellTwo\)-saturated.
\end{corollary}

\begin{proof}
The first assertion is the contrapositive of \cref{thm:corrected}; the second adds the hypotheses of Godefroy's problem.
\end{proof}

\begin{remark}
The phrase ``unless \(X\) is Hilbertian'' in the submitted theorem is unnecessary.  Every infinite-dimensional separable Hilbert space is isomorphic to \(\ellTwo\), hence is \(\ellTwo\)-saturated and does not satisfy the hypothesis of \cref{thm:corrected}.
\end{remark}

\begin{center}
\fcolorbox{darkgreen}{softgreen}{\begin{minipage}{0.90\textwidth}
\textbf{The repaired argument in one line:}
\[
  X\text{ not }\ellTwo\text{-saturated}
  \Longrightarrow
  \exists W\subseteq X\text{ with no }\ellTwo
  \overset{\text{Anisca}}{\Longrightarrow}
  |\Spec(W)|=\infty.
\]
\end{minipage}}
\end{center}

\section{Literature status: no full answer found}

The 2010 source survey presents the question as open and notes only partial results \cite{DLT2010}.  The most direct modern status statement located is in Cuellar Carrera--de Rancourt--Ferenczi (2023): even Johnson's special question asking for a separable Banach space with exactly two subspace isomorphism types remains open, and it is not known whether a separable non-Hilbertian space can have at most countably many such types \cite[pp.~3539--3540]{CCRF2023}.

A still stronger conjecture of Ferenczi--Rosendal asserts that every separable non-Hilbertian Banach space is \emph{ergodic}, meaning that \(\Ezero\) Borel-reduces to isomorphism between its subspaces \cite{FRergodic}.  Ergodicity gives continuum many pairwise non-isomorphic subspaces.  De Rancourt--Kurka (2026) explicitly describe this conjecture as still widely open, while proving it for every non-Hilbertian Orlicz sequence space and for the Kalton--Peck family \(\ell_2(\varphi)\), including \(Z_2\) \cite{dRK2026}.

Accordingly, the search did not trigger the requested ``stop because the literature already contains a full answer'' condition.  No proof or counterexample to Godefroy's problem was found.  The conclusion is necessarily a literature-status assessment rather than a claim about unpublished work.

\section{What a counterexample would now have to look like}

Assume that \(X\) is a counterexample.  Then \(\Spec(X)\) is finite, so every subspace of \(X\) is non-ergodic.  Combining \cref{cor:saturated} with modern structure theorems gives the following restrictive profile.

\begin{proposition}[Necessary features of a counterexample]\label{prop:profile}
A counterexample \(X\) must satisfy all of the following.
\begin{enumerate}
  \item \(X\) is \(\ellTwo\)-saturated.
  \item \(X\) is near Hilbert: it has type \(p\) and cotype \(q\) for every \(p<2<q\) \cite{Cuellar2018}.
  \item No non-Hilbertian subspace of \(X\) is asymptotically Hilbertian, since such a subspace would be ergodic by Anisca's theorem \cite{Anisca2010}.
  \item Every subspace of \(X\) contains a further subspace isomorphic to its square; this follows from the Ferenczi--Rosendal dichotomy because no subspace of \(X\) can contain continuum many non-isomorphic subspaces \cite[Corollary 14]{FRnumber}.
  \item \(X\) contains a non-Hilbertian \(d_2\)-minimal subspace, i.e. a non-Hilbertian subspace that embeds into each of its non-Hilbertian subspaces \cite[Theorem 6.5]{CCRF2023}.
  \item After passing to a further non-Hilbertian subspace, one reaches either a space with an unconditional basis or a space that is simultaneously \(d_2\)-minimal and \(d_2\)-HI (it contains no direct sum of two non-Hilbertian subspaces) \cite[Theorem 6.23]{CCRF2023}.
  \item It cannot be a non-Hilbertian Orlicz sequence space, the Kalton--Peck space \(Z_2\), or any non-Hilbertian member of the family \(\ell_2(\varphi)\) treated in \cite{dRK2026}.
\end{enumerate}
\end{proposition}

The fourth item is especially significant.  The 2023 paper reports no known example of a nontrivial \(d_2\)-minimal space and asks whether one exists at all \cite[Question 6.8]{CCRF2023}.  It also conjectures that a Banach space cannot be simultaneously \(d_2\)-minimal and \(d_2\)-HI \cite[Conjecture 6.26]{CCRF2023}.  On the other hand, Argyros--Manoussakis--Motakis constructed an \(\ellTwo\)-saturated \(d_2\)-HI space \cite{AMM2025}; thus \(d_2\)-HI behavior itself is possible, but the additional minimality demanded by a non-ergodic candidate remains the hard point.

For the exact-two-types case (a ``Johnson space''), still more is known: such a space would be \(d_2\)-minimal, would have a Schauder basis, and would have an unconditional basis exactly when it is isomorphic to its square \cite[Proposition 6.7 and Corollary 6.16]{CCRF2023}.  Even this sharply constrained special case is unresolved.

\section{Attempts to push beyond the corrected theorem}

This section records the principal proof routes considered and the precise obstruction met by each.  These are not merely stylistic failures: each route runs into a known open structural issue.

\subsection{Trying to repair the minimal-class argument}

The most direct repair would be a Schroeder--Bernstein theorem inside the class of subspaces under consideration.  Mutual embeddability would then become isomorphism, and a minimal representative would be homogeneous.  In general this is false by Gowers' counterexample \cite{GowersSB}.  Standard positive substitutes, such as Pe\l czy\'nski's decomposition method, require strong complemented-decomposition and square-isomorphism hypotheses that are not supplied by finite subspace spectrum.

Even a restricted Schroeder--Bernstein statement for \(\ellTwo\)-saturated, finite-spectrum spaces would be a substantial new theorem.  Moreover, it would still have to control non-Hilbertian subspaces relative to the ever-present \(\ellTwo\) class; ordinary minimality is too weak, and the relevant replacement is \(d_2\)-minimality.

\subsection{A finite Borel coloring and Ramsey homogenization}

Because isomorphism between separable Banach spaces is analytic, the individual isomorphism classes inside \(\Sub(X)\) become Borel when there are only finitely many of them: the complement of one class is then a finite union of analytic classes.  It is tempting to apply a Ramsey theorem and obtain a subspace all of whose further subspaces have one color, which would be homogeneous and therefore Hilbertian.

The missing ingredient is exactly that available infinite-dimensional Ramsey theorems do not homogenize arbitrary Borel colorings of \emph{all} closed subspaces under inclusion.  On the coding spaces of subsequences and block sequences they give strong conclusions, but those conclusions concern coordinate or block subspaces, or strategic largeness, rather than every closed subspace.  The local dichotomies of \cite{CCRF2023} can be viewed as the refined outcome of this strategy: they produce \(d_2\)-minimality or tightness, not full homogeneity.  Thus finite Borel complexity alone does not repair the submitted proof.

\subsection{Stabilizing subsequence or block subspaces}

Every infinite-dimensional Banach space contains a basic sequence, so one can encode coordinate or block subspaces by infinite subsets of \(\N\).  Baire-category and Ramsey arguments of Ferenczi--Rosendal give a sharp dichotomy for spaces with unconditional bases: either one obtains a perfect family of pairwise non-isomorphic coordinate subspaces, or strong self-similarity follows \cite{FRnumber}.  In the non-ergodic case, a space with an unconditional basis is isomorphic to its square and to many coordinate extensions \cite{FRnumber,CCRF2023}.

The obstacle is that coordinate and block subspaces are only a small subclass of all subspaces.  Stabilizing their isomorphism types does not make the ambient space homogeneous.  Pushing this line toward all subspaces collides with unresolved block-homogeneity questions; the source survey's Problem 6 is one manifestation of precisely this gap \cite{DLT2010}.

\subsection{Square-saturation and decomposition methods}

Ferenczi--Rosendal proved that every Banach space either has continuum many pairwise non-isomorphic subspaces or is saturated with subspaces isomorphic to their squares \cite[Corollary 14]{FRnumber}.  Hence a finite-spectrum candidate has abundant square-isomorphic subspaces.  This does not finish the problem.  In a \(d_2\)-HI space a non-Hilbertian square-isomorphic subspace is impossible, but the square subspaces supplied by the theorem may all be Hilbertian; \(\ellTwo\)-saturation already provides those.  In the unconditional branch, square-isomorphism supports Pe\l czy\'nski-type decomposition arguments, but one still lacks complemented minimality or complemented mutual embeddings, so cancellation cannot be justified.

\subsection{Coding isomorphism types with finite-dimensional distortion}

A common way to build many non-isomorphic subspaces is to arrange finite-dimensional pieces with different Euclidean distortion and encode a set \(A\subseteq\N\) by selecting pieces indexed by \(A\).  This succeeds in asymptotically Hilbertian and Orlicz settings, and underlies several ergodicity results \cite{Anisca2010,dRK2026}.

A hypothetical counterexample must avoid every non-Hilbertian asymptotically Hilbertian subspace.  Worse, in a \(d_2\)-minimal space, highly non-Euclidean finite-dimensional pieces occur uniformly inside every non-Hilbertian subspace \cite[Lemma 6.14]{CCRF2023}.  Naive distortion profiles therefore recur everywhere and cease to distinguish subspaces.  In the \(d_2\)-HI branch, direct-sum coding is additionally blocked because no two non-Hilbertian pieces may form a direct sum.

\subsection{Eliminating the \texorpdfstring{\(d_2\)}{d2}-HI branch}

Theorem 6.23 of \cite{CCRF2023} reduces a non-ergodic non-Hilbertian space to an unconditional branch or a simultaneous \(d_2\)-minimal/\(d_2\)-HI branch.  Proving Conjecture 6.26 of that paper would eliminate the latter.  The 2025 construction of an \(\ellTwo\)-saturated \(d_2\)-HI space shows that one cannot simply prove that \(\ellTwo\)-saturation and \(d_2\)-HI are incompatible \cite{AMM2025}; one must use minimality in an essential way.

A plausible target is to show that sufficiently deep self-embeddings of a \(d_2\)-HI space necessarily leave a non-Hilbertian complement or create two separated non-Hilbertian subspaces.  Existing results only establish such complementability conclusions for specific examples, not for all \(d_2\)-HI spaces \cite[Section 6.2]{CCRF2023}.

\subsection{The twisted Hilbert route}

Counterexamples must be near Hilbert, so twisted Hilbert spaces are natural candidates.  The latest broad result rules out \(Z_2\), every \(\ell_2(\varphi)\) in the Kalton--Peck family, and twisted Hilbert spaces generated by interpolation between Orlicz sequence spaces: each is Hilbertian or ergodic \cite{dRK2026}.  The paper leaves open whether every non-Hilbertian separable twisted Hilbert space is ergodic.  Thus this route has narrowed but not closed.

\section{A precise remaining target}

The corrected theorem reduces the original question to the following form.

\begin{question}[Core finite-spectrum problem]\label{q:core}
Can a separable, non-Hilbertian, \(\ellTwo\)-saturated Banach space have only finitely many isomorphism classes of closed infinite-dimensional subspaces?
\end{question}

A negative answer to Question~\ref{q:core} is exactly an affirmative answer to Godefroy's Problem 7.  The modern reductions suggest two more focused targets:
\begin{enumerate}
  \item prove that no nontrivial \(d_2\)-minimal Banach space has finite subspace spectrum;
  \item or, more ambitiously, prove that every nontrivial \(d_2\)-minimal separable Banach space is ergodic.
\end{enumerate}
The second statement is a special case of the Ferenczi--Rosendal conjecture and would yield continuum many, not merely countably many, non-isomorphic subspaces.

\section{Conclusion}

The submitted partial result has the correct conclusion but an invalid proof.  Minimality in the embeddability preorder yields mutual embeddability, and the Banach-space Schroeder--Bernstein theorem is false.  The clean repair is the finite-spectrum theorem of Anisca: a non-\(\ellTwo\)-saturated space contains a subspace with no \(\ellTwo\), which therefore cannot have finite subspace spectrum.

No full solution or counterexample was found in the literature through 20 June 2026.  Current results force any counterexample into a remarkably narrow class: finite-spectrum, \(\ellTwo\)-saturated, near Hilbert, non-asymptotically-Hilbertian, and containing a nontrivial \(d_2\)-minimal core.  The remaining difficulty is not the corrected partial theorem but the unresolved structure of such \(d_2\)-minimal spaces, particularly in the unconditional and \(d_2\)-HI branches.

\begin{thebibliography}{99}

\bibitem{AMM2025}
S. A. Argyros, A. Manoussakis, and P. Motakis,
\emph{Variants of the James Tree space},
Israel J. Math. 265 (2025), 1--77; arXiv:2012.06286.

\bibitem{Anisca2007}
R. Anisca,
\emph{On the structure of Banach spaces with an unconditional basic sequence},
Studia Math. 182 (2007), 67--85.

\bibitem{Anisca2010}
R. Anisca,
\emph{The ergodicity of weak Hilbert spaces},
Proc. Amer. Math. Soc. 138 (2010), 1405--1413.

\bibitem{Cuellar2018}
W. Cuellar Carrera,
\emph{Non-ergodic Banach spaces are near Hilbert},
Trans. Amer. Math. Soc. 370 (2018), 8691--8707.

\bibitem{CCRF2023}
W. Cuellar Carrera, N. de Rancourt, and V. Ferenczi,
\emph{Local Banach-space dichotomies and ergodic spaces},
J. Eur. Math. Soc. 25 (2023), 3537--3598.

\bibitem{dRK2026}
N. de Rancourt and O. Kurka,
\emph{The ergodicity of Orlicz sequence spaces},
J. Funct. Anal. 290 (2026), no. 5, Article 111270;
arXiv:2501.17756v3.

\bibitem{DLT2010}
P. Dodos, J. Lopez-Abad, and S. Todorcevic,
\emph{Banach spaces and Ramsey Theory: some open problems},
arXiv:1006.2668v1 (2010).

\bibitem{FRergodic}
V. Ferenczi and C. Rosendal,
\emph{Ergodic Banach spaces},
Adv. Math. 195 (2005), 259--282.

\bibitem{FRnumber}
V. Ferenczi and C. Rosendal,
\emph{On the number of non-isomorphic subspaces of a Banach space},
Studia Math. 168 (2005), 203--216.

\bibitem{GowersSB}
W. T. Gowers,
\emph{A solution to the Schroeder--Bernstein problem for Banach spaces},
Bull. London Math. Soc. 28 (1996), 297--304.

\end{thebibliography}

\end{document}
